% Ad Familiares 7.6 --- headnote
% ad-familiares-07-06, May 54 BC, Cumanum or Pompeianum

\textit{Cicero to Trebatius, written from the Cuman or
Pompeian villa in May 54 BC. The third letter of the Trebatius
sequence after the recommendation to Caesar (\textit{Fam.}
7.5, April 54) and the parallel commendation to Quintus
(implicit in \textit{Q. fr.} 2.12, May 54). Trebatius is now
with Caesar's army, on the way to or in Britain, and is
homesick.}

\textit{The body is a sustained literary play with Ennius's
\textit{Medea}, the Roman tragedy adapted from Euripides. The
silly notions and longings for the city are to be set aside
and the campaign endured: we your friends will excuse you,
as the Corinthian matrons excused Medea, who, ``with hands
most thoroughly whitened with chalk'' (a touch of female
ritual --- the chalked palms of formal courtesy), persuaded
them not to think the worse of her for being away from her
country. The two verse quotations from Ennius are the load:
``\textit{the wealthy matrons, the gentlewomen, that held
Corinth's high citadel}'' (the matrons addressed in Medea's
defence) and the apothegm ``\textit{many have managed both
their own affair, and the commonwealth, well, far from the
fatherland; many, who would spend their life at home, on that
very account have been disapproved.}''}

\textit{The closing line returns to the same play, with the
double tag of Roman home wisdom: ``take care that in Britain
you are not tricked by the chariot-fighters'' --- the British
\textit{essedarii}, the war-chariot drivers Caesar's
\textit{De Bello Gallico} would describe; and Ennius again,
``\textit{he who, being wise, cannot help his own self, is
wise to no purpose}.''  The combination is typical Ciceronian
\textit{urbanitas}: the educated joke at Trebatius's expense
that takes its weight from the great old tragic poet
Trebatius would have read with him.}
